\section{Theoretical Results}\label{sec:results-theory}

This section reports what the model says before any race is consulted. All
numbers are produced by the repository's \texttt{run\_all} script from the
parameter registry; the reference swimmer is calibrated so that the optimal
200-yard time is exactly 100.0~s (1:40.0), which makes penalties readable as
seconds lost per 100~s of racing.

\subsection{The optimum and the cost of missing it}

Under the pure budget (M0) the optimum is four laps of 25.000~s, and the
closed form, the split-space SQP solve and the direct ODE integration agree
on it to $10^{-11}$~s, as Theorem~\ref{thm:even} requires. The interesting
quantity is not the optimum but the penalty for missing it, because the
penalty is what a coach would act on. Figure~\ref{fig:penalty} plots the
race time when the opening lap is forced to a fixed error and the remaining
three laps are re-optimized around it. The penalty is quadratic near the
optimum and small: a 2\% opening error costs 0.015~s, 5\% costs 0.09~s, and
10\% costs 0.36~s (too fast) or 0.39~s (too slow). Named strategies at the
same energy budget tell the same story (Figure~\ref{fig:strategies}): a
smooth ramp spanning 1.5 points from first lap to last, in either
direction, costs 0.055~s; the ``aggressive opening'' profile, 7.4\% fast on
lap~1, costs 0.27~s; a ``very aggressive'' 13.0\% opening costs 0.94~s; a
conservative 6.4\% slow opening costs 0.18~s. The corrected elite 200-m finalist profile of
\citet{robertson2009analysis} costs 0.009~s against the even optimum, and
the mean corrected profile of our own 80-race first pilot costs 0.075~s. A
Monte Carlo over the physiological block (240 draws with $R$, $E_0$, $C_D$,
$A$, $\eta_p$, $\eta_g$ perturbed log-normally at 10\% coefficient of
variation; Figure~\ref{fig:robust}) leaves this ordering intact in every
draw, with the very-aggressive penalty spanning 0.83--1.06~s and the smooth
ramps 0.049--0.063~s across the 5th--95th percentiles. The ranking of
pacing errors is therefore robust to parameter ignorance of the size we
actually have; the magnitude is robust to within about 15\%.

\begin{figure}[htbp]
\centering
\includegraphics[width=0.62\linewidth]{fig05_opening_penalty}
\caption{Race-time penalty for a forced opening-lap error, remaining laps
re-optimized (reference swimmer, M0). The penalty is quadratic near the
optimum; a 5\% opening error costs about 0.09~s of a 100~s race.}
\label{fig:penalty}
\end{figure}

\begin{figure}[htbp]
\centering
\includegraphics[width=0.74\linewidth]{fig03_strategy_comparison}
\caption{Named pacing strategies at matched energy budget. The observed
elite profile, once the dive credit is added back, lies within 0.01~s of the
even optimum.}
\label{fig:strategies}
\end{figure}

\begin{figure}[htbp]
\centering
\includegraphics[width=0.7\linewidth]{fig15_penalty_robustness}
\caption{Pacing penalties under 240 seeded Monte Carlo draws of the
physiological block (10\% CV). The ordering of strategies survives every
draw; magnitudes move by about 15\%.}
\label{fig:robust}
\end{figure}

\subsection{The engine sets the speed, not the shape}

Theorem~\ref{thm:shape} says the optimal shape is independent of every
engine and hull parameter; Figure~\ref{fig:elasticity} shows how strongly
those parameters act on everything else. A 10\% change in either the drag
area $A$ or the drag coefficient $C_D$ moves the optimal race time by
7.4~s (elasticity 0.37), the same change in either efficiency moves it by
7.4~s the other way, and $R$ and $E_0$ move it by 5.3~s and 2.1~s
(elasticities 0.26 and 0.11). Not one of these changes moves the first-lap
share by more than $10^{-8}$. Only the fatigue couplings and the kinetic
time constant do anything to shape or to timing beyond the budget: stepping
$\beta_x$ from 0 to 0.05 lowers the optimal first-lap share by 0.15 points
and costs 0.91~s, while stepping $\tau$ from 0 to 2~s costs 0.53~s and
leaves the shape exactly even, as Proposition~\ref{prop:kinetics} says it
must. The archetype experiment (Figure~\ref{fig:archetypes}) makes the same
point at the level a coach would recognize: seven swimmers calibrated to the
same 1:43.0 --- aerobic-dominant, anaerobic-dominant, high-drag, low-drag
and reference --- all receive an identical, exactly even prescription
(shape spread below $10^{-8}$), and only the two fatigue archetypes differ,
the fatigue-prone swimmer ($\beta_x=0.35$) at $(.2405,.2470,.2533,.2592)$
and the fatigue-resistant one ($\beta_x=0.08$) at $(.2476,.2492,.2508,.2524)$.
The intuition that a stronger athlete has earned a faster opening lap finds
no support here; what would earn one is a \emph{smaller} fatigue coupling.

\begin{figure}[htbp]
\centering
\includegraphics[width=0.7\linewidth]{fig13_elasticity}
\caption{Race-time elasticities under $\pm10\%$ parameter perturbations. The
drag and efficiency block dominates race time; none of the engine or hull
parameters moves the optimal shape.}
\label{fig:elasticity}
\end{figure}

\begin{figure}[htbp]
\centering
\includegraphics[width=0.74\linewidth]{fig12_archetypes}
\caption{Seven archetypes calibrated to the same target time. Only the
fatigue couplings change the prescribed shape.}
\label{fig:archetypes}
\end{figure}

\subsection{Mechanisms as shapes}

Figure~\ref{fig:mechanisms} stacks the mechanisms on the reference swimmer.
Oxygen kinetics with $\tau=20$~s costs 5.5~s and with $\tau=30$~s costs
8.4~s, in both cases with an exactly even optimum. Position-coupled decay at
$\beta_x=0.15$ and $0.30$ yields near-linear positive splits,
$(.2456,.2486,.2515,.2543)$ and $(.2417,.2474,.2528,.2581)$; reserve-coupled
decay at $\beta_E=0.30$ yields a strong negative split
$(.2815,.2605,.2382,.2198)$. The heat-map of the first-lap share over the
$(\beta_x,p)$ plane (Figure~\ref{fig:heatmap}) shows the whole M3 family at
once: the share falls monotonically in $\beta_x$ and rises toward $0.25$ as
the cost exponent $p$ grows, exactly the $w_i^{1/p}$ dependence of
Eq.~\eqref{eq:shape}. Two consequences matter for the empirical section.
First, the differences between the three shape classes are of order one to
three percentage points per lap, an order of magnitude larger than the
race-to-race scatter of shares within a class, so even a few hundred races
can separate the classes if the start is handled. Second, the differences
\emph{within} the positive class, M3 against M4, are of order a tenth of a
point per lap at the registry couplings, which is the same order as the
effect of moving the dive credit by a few tenths of a second. The theory
therefore predicts, before any data, where the empirical comparison will be
decisive and where it will be hostage to the start.

\begin{figure}[htbp]
\centering
\includegraphics[width=0.74\linewidth]{fig09_theory_vs_observed}
\caption{Model optima under each mechanism against observed strategy shapes.
The mechanisms collapse onto three shape classes: even, negative, positive.}
\label{fig:mechanisms}
\end{figure}

\begin{figure}[htbp]
\centering
\includegraphics[width=0.55\linewidth]{fig14_shape_heatmap}
\caption{Optimal first-lap share over the position-fatigue coefficient and
the cost exponent, the only two parameters that enter the M3 shape.}
\label{fig:heatmap}
\end{figure}

\subsection{What the working swimmer prescribes}

The working swimmer used for prescriptions is the M3 variant at the registry
coupling $\beta_x=0.28$, recalibrated to 1:40.0. Its optimum is
$(.2422,.2475,.2527,.2576)$: 24.22, 24.75, 25.27, 25.76~s. Against it, even
pacing costs 0.056~s, a negative split of the same magnitude costs 0.22~s,
the conservative opening 0.43~s and the aggressive opening 0.08~s. Which
direction of error is cheaper thus depends on which mechanism is true, and
that is the question the data are asked next.

\FloatBarrier
